\scanpage{160}
\mspage{109}

telle que $p\in\unsure{W|\mathcal W}$, la flèche $g:X'\mathbin{-}X''$ qui s'en \unsure{déduit} par cochangement de base à partir de l'objet $X$ de $U\backslash\mathrm{Cat}$, est aussi dans $W$.

\uline{Dém. du corollaire} : on applique le th. 3 à l'immersion ouverte $U'\overset{i'}{\hookrightarrow}X'$ \unsure{déduite} de $U\overset{i}{\hookrightarrow}X$ par le cochangement de base $U\xrightarrow{p}U'$, en notant de plus que l'hypothèse du th. 3 sur $i$ est \unsure{préservée} en passant de $i$ à $i'$.

Il faut \unsure{quand} même introduire une terminologie adéquate :

\uline{Définition 2} : a) Une flèche $f:X\mathbin{-}Y$ est dite \uline{strict\deleted{\unsure{em}}ement $W$-asph.} si pour tout $y\in\operatorname{Ob}Y$, $f^{*}(y)=(x\longmapsto\operatorname{Hom}_{Y}(f(x),y))$ est un objet \unsure{tot.} $W$-asph. de $X^{\wedge}$ (cf déf. 1 page 88). Elle est dite \deleted{strictement} \uline{totalement} $W$-\uline{asphérique} si, pour tout $y\in\operatorname{Ob}Y$, $\mathrm{cat.}\ X_{/y}\ (\simeq X_{/f^{*}(y)})$ est \deleted{\unsure{t}}un objet \unsure{tot.} $W$-asphérique \unsure{de|ds}.

b) Une immersion ouverte $i:U\mathbin{-}X$ est dite \deleted{\unsure{(faiblement)}} \inserted{$W$-}dwyérienne, \inserted{(resp. \deleted{totalement} dwyérienne)} \inserted{si} l'immersion \unsure{induite} $i_0:U\mathbin{-}\bar U$ est \deleted{\unsure{strict}}strictement \inserted{(resp. totalement)} \uline{$W$-asphérique}. Une immersion fermée est \deleted{\unsure{dite ou non}} \inserted{$j:J\mathbin{-}X$} est dite \inserted{$W$-}dwyérienne (resp. tot. $W$-dwyérienne) ssi $j^{\mathrm o}:\unsure{I|J}^{\mathrm o}\mathbin{-}X^{\mathrm o}$ est une imm. ouverte $W$-dwyérienne (resp. tot. $W$-dwyérienne).

c) Une \deleted{\unsure{immersion ouverte ou fermée est}} \inserted{flèche $i:U\mathbin{-}X$ \unsure{ds|de} $\mathrm{Cat}$ (pas néc. une imm.)}
% archive sheet number: 160
